% Concepts from practice --- what an exercise set taught, built from session records rather than from
% the solutions a second time. One per set, beside the solutions.
\documentclass[a4paper,10pt]{article}
\IfFileExists{preamble.tex}{\input{preamble.tex}}{\IfFileExists{../../../../../60 Templates/preamble.tex}{\input{../../../../../60 Templates/preamble.tex}}{\IfFileExists{../../../../60 Templates/preamble.tex}{\input{../../../../60 Templates/preamble.tex}}{\input{../../../60 Templates/preamble.tex}}}}
\begin{document}
\ArtifactHeader{\ProjectCode}{<meeting and existing Research-map key or set>}{<the exercise set's name> --- Concepts}{<YYYY-MM-DD>}%
  {<the assigned source and the session records this was built from>}

\section{What the exercises were for}
<The one thing the exercise set was built to make land, in a sentence. If that sentence is hard to write,
the records are the place to go back to, not the solutions.>

\section{<The idea the exercise set develops>}

\begin{defn}[<term>]
<Restated in the form the questions actually used it, which is often narrower than the source's.>
\end{defn}

\begin{remark}
<The move the exercise set kept asking for, named so it is recognisable next time it is needed.>
\end{remark}

\begin{warning}
<Where the working went wrong, taken from the records rather than guessed --- this is the part a
later session reads first.>
\end{warning}

\section{Where this shows up again}
\begin{itemize}
  \item <A question elsewhere in the project that turns on the same idea, named by pointer.>
  \item <The learning unit whose material this set exercised, by Research-map key.>
\end{itemize}

\begin{checkyourself}
  \item <The question from the set worth being able to do cold.>
\end{checkyourself}
\end{document}
